% 普莱费尔公理（综述）
% license CCBYSA3
% type Wiki

本文根据 CC-BY-SA 协议转载翻译自维基百科\href{https://en.wikipedia.org/wiki/Playfair\%27s_axiom}{相关文章}

\begin{figure}[ht]
\centering
\includegraphics[width=14.25cm]{./figures/047cd41d4213cc97.png}
\caption{} \label{fig_PLFRgl_1}
\end{figure}
在几何学中，普莱费尔公理是一个可以替代欧几里得第五公设（平行公设）的公理：

在一个平面内，给定一条直线和一个不在该直线上的点，至多只能作一条经过该点且与已知直线平行的直线。[1]

它在欧几里得几何的语境下与**欧几里得平行公设**等价\[2]，并以苏格兰数学家约翰·普莱费尔（John Playfair）的名字命名。这里的“**至多一条**”这一限定就足够了，因为可以从前四条公理推导出：给定一条直线 $L$ 和一个不在 $L$ 上的点 $P$，至少存在一条平行线。推导过程如下：

* **作垂线**：利用公理和已建立的定理，可以作出一条通过点 $P$ 的直线，并且垂直于直线 $L$。
* **作另一条垂线**：再作一条通过点 $P$ 的直线，使其垂直于第一条垂线。
* **得到平行线**：由于平行线的定义（即根据第四条公理，内错角相等），这条第二条垂线与 $L$ 平行。

该命题通常写作“存在唯一一条平行线”。在欧几里得的《几何原本》中，两条直线被称为平行，当且仅当它们永不相交；而其他对平行线的刻画并未被使用。\[3]\[4]

这一公理不仅用于欧几里得几何，也用于更广义的**仿射几何**研究中，在那里“平行性”的概念居于核心地位。在仿射几何的背景下，需要使用普莱费尔公理的更强形式（即将“至多一条”替换为“唯一一条”），因为中性几何的公理体系缺乏提供存在性的证明。普莱费尔版的公理已如此广泛流行，以至于它常被称为**欧几里得的平行公理**\[5]，尽管这并不是欧几里得本人提出的版本。
